\textbf{Phần 1. Trắc nghiệm lựa chọn}
\setcounterpageref{ex}{0}

\begin{ex} %[2D8N44]
Kết quả của phép tính $15 : 5$ là
\choice
{$2$}
{$4$}
{\True $3$}
{$5$}
\loigiai{
$15 : 5 = 3$.\\
\textbf{Đáp án của câu hỏi là C}
}
\end{ex}

\begin{ex} %[2D8N44]
Kết quả của phép tính $25 : 5$ là
\choice
{$4$}
{$6$}
{$7$}
{\True $5$}
\loigiai{
$25 : 5 = 5$.\\
\textbf{Đáp án của câu hỏi là D}
}
\end{ex}

\begin{ex} %[2D8H44]
Phép tính nào sau đây có kết quả bằng 8?
\choice
{\True $40 : 5$}
{$35 : 5$}
{$45 : 5$}
{$30 : 5$}
\loigiai{
$40 : 5 = 8$.\\
\textbf{Đáp án của câu hỏi là A}
}
\end{ex}

\begin{ex} %[2D8H44]
Phép tính nào trong bảng chia 5 có kết quả lớn nhất?
\choice
{$45 : 5$}
{\True $50 : 5$}
{$40 : 5$}
{$35 : 5$}
\loigiai{
$50 : 5 = 10$ (lớn nhất).\\
\textbf{Đáp án của câu hỏi là B}
}
\end{ex}

\begin{ex} %[2D8H44]
Trong bảng chia 5, các số bị chia đều tận cùng là chữ số nào?
\choice
{$1$ hoặc $5$}
{$2$ hoặc $5$}
{\True $0$ hoặc $5$}
{Chỉ chữ số $0$}
\loigiai{
Tận cùng là 0 hoặc 5.\\
\textbf{Đáp án của câu hỏi là C}
}
\end{ex}

\begin{ex} %[2D8H44]
Điền số thích hợp: $35 : 5 = \dots$
\choice
{$6$}
{$8$}
{$9$}
{\True $7$}
\loigiai{
$35 : 5 = 7$.\\
\textbf{Đáp án của câu hỏi là D}
}
\end{ex}

\begin{ex} %[2D8V44]
Có 20 quyển vở chia đều cho 5 bạn. Mỗi bạn nhận được số quyển vở là
\choice
{$3$ quyển}
{\True $4$ quyển}
{$5$ quyển}
{$6$ quyển}
\loigiai{
$20 : 5 = 4$ quyển vở.\\
\textbf{Đáp án của câu hỏi là B}
}
\end{ex}

\begin{ex} %[2D8V44]
Một bao đường nặng $30\text{ kg}$ chia đều vào 5 túi. Mỗi túi nặng số ki-lô-gam là
\choice
{\True $6\text{ kg}$}
{$5\text{ kg}$}
{$7\text{ kg}$}
{$8\text{ kg}$}
\loigiai{
$30 : 5 = 6\text{ kg}$.\\
\textbf{Đáp án của câu hỏi là A}
}
\end{ex}

\begin{ex} %[2D8V44]
Tính giá trị biểu thức: $45 : 5 + 28$.
\choice
{$35$}
{$39$}
{$41$}
{\True $37$}
\loigiai{
$9 + 28 = 37$.\\
\textbf{Đáp án của câu hỏi là D}
}
\end{ex}

\begin{ex} %[2D8V44]
Điền dấu thích hợp: $40 : 5 \dots 18 : 2$.
\choice
{\True $<$}
{$>$}
{$=$}
{Không so sánh được}
\loigiai{
$40 : 5 = 8 < 18 : 2 = 9$.\\
\textbf{Đáp án của câu hỏi là A}
}
\end{ex}

\begin{ex} %[2D8V44]
Có 50 lít dầu rót đều vào 5 thùng. Mỗi thùng có số lít dầu là
\choice
{$8\text{ l}$}
{$12\text{ l}$}
{\True $10\text{ l}$}
{$15\text{ l}$}
\loigiai{
$50 : 5 = 10\text{ l}$.\\
\textbf{Đáp án của câu hỏi là C}
}
\end{ex}

\begin{ex} %[2D8C44]
Có 30 bông hoa cắm đều vào 5 bình hoa. Sau đó cắm thêm vào mỗi bình 2 bông hoa nữa. Hỏi mỗi bình có tất cả bao nhiêu bông hoa?
\choice
{$7$ bông}
{\True $8$ bông}
{$9$ bông}
{$10$ bông}
\loigiai{
Ban đầu mỗi bình: $30 : 5 = 6$ bông.\\
Lúc sau mỗi bình: $6 + 2 = 8$ bông hoa.\\
\textbf{Đáp án của câu hỏi là B}
}
\end{ex}

\textbf{Phần 2. Trắc nghiệm đúng sai}
\setcounterpageref{ex}{0}

\begin{ex} %[2D8H44]
Xét tính đúng sai của các phép tính trong bảng chia 5:
\choiceTF
{\True $10 : 5 = 2$.}
{\True $25 : 5 = 5$.}
{$45 : 5 = 8$.}
{\True $45 : 5 = 9$.}
\loigiai{
\begin{itemize}
    \item $10 : 5 = 2$: Đúng.
    \item $25 : 5 = 5$: Đúng.
    \item $45 : 5 = 9 \neq 8$. Ý 3 Sai, Ý 4 Đúng.
\end{itemize}
\begin{itemize}
    \item a) \textbf{Đúng}.
    \item b) \textbf{Đúng}.
    \item c) \textbf{Sai}.
    \item d) \textbf{Đúng}.
\end{itemize}
}
\end{ex}

\begin{ex} %[2D8V44]
Có 35 quả cam chia đều vào 5 đĩa. Xét tính đúng sai:
\choiceTF
{\True Mỗi đĩa có 7 quả cam ($35 : 5 = 7$).}
{\True 3 đĩa như thế có 21 quả cam ($7 \times 3 = 21$).}
{Nếu lấy bớt 1 đĩa thì còn lại 30 quả cam.}
{\True Khi đó còn lại $7 \times 4 = 28$ quả cam.}
\loigiai{
\begin{itemize}
    \item Mỗi đĩa 7 quả: Đúng.
    \item 3 đĩa có 21 quả: Đúng.
    \item Còn lại: $7 \times 4 = 28 \neq 30$. Ý 3 Sai, Ý 4 Đúng.
\end{itemize}
\begin{itemize}
    \item a) \textbf{Đúng}.
    \item b) \textbf{Đúng}.
    \item c) \textbf{Sai}.
    \item d) \textbf{Đúng}.
\end{itemize}
}
\end{ex}

\textbf{Phần 3. Trả lời ngắn}
\setcounterpageref{ex}{0}

\begin{ex} %[2D8H44]
Tính: $20 : 5$.
\shortans{4}
\loigiai{
\textbf{Đáp số:} $4$\\
$20 : 5 = 4$.
}
\end{ex}

\begin{ex} %[2D8H44]
Tính: $35 : 5$.
\shortans{7}
\loigiai{
\textbf{Đáp số:} $7$\\
$35 : 5 = 7$.
}
\end{ex}

\begin{ex} %[2D8V44]
Tính: $40 : 5$.
\shortans{8}
\loigiai{
\textbf{Đáp số:} $8$\\
$40 : 5 = 8$.
}
\end{ex}

\begin{ex} %[2D8V44]
Tính giá trị biểu thức: $50 : 5 - 6$.
\shortans{4}
\loigiai{
\textbf{Đáp số:} $4$\\
$10 - 6 = 4$.
}
\end{ex}

\begin{ex} %[2D8V44]
Có 45 cái kẹo chia đều cho 5 bạn. Hỏi mỗi bạn được bao nhiêu cái kẹo?
\shortans{9}
\loigiai{
\textbf{Đáp số:} $9$\\
$45 : 5 = 9$ cái kẹo.
}
\end{ex}

\begin{ex} %[2D8C44]
Một lớp có 35 học sinh chia đều thành 5 tổ. Tổ Một có 3 bạn nữ. Hỏi tổ Một có bao nhiêu bạn nam?
\shortans{4}
\loigiai{
\textbf{Đáp số:} $4$\\
Mỗi tổ có: $35 : 5 = 7$ bạn.\\
Số bạn nam: $7 - 3 = 4$ bạn.
}
\end{ex}
